\documentclass[
    14pt,
    candidate,
    times
]{disser}

% Кодировка и язык
\usepackage[T2A]{fontenc} % поддержка кириллицы
\usepackage[utf8]{inputenc} % кодировка исходного текста
\usepackage[english,russian]{babel} % переключение языков

% Геометрия страницы и графика
\usepackage[left=2.5cm, right=2cm, top=2cm, bottom=2cm]{geometry} % поля страницы
\usepackage{graphicx} % подключение графики
\usepackage{pdfpages} % вставка pdf-страниц

% Таблицы
\usepackage{array} % расширенные возможности для работы с таблицами
\usepackage{tabularx} % автоматический подбор ширины столбцов
\usepackage{dcolumn} % выравнивание чисел по разделителю

% Математика
\usepackage{bm} % полужирное начертание для математических символов
\usepackage{amsmath} % дополнительные математические возможности
\usepackage{amssymb} % дополнительные математические символы

% Библиография и ссылки
\usepackage{cite} % поддержка цитирования
\usepackage[hyperfootnotes=false,hidelinks]{hyperref} % создание гиперссылок
\providecommand{\selectlanguageifdefined}[1]{%
  \expandafter\ifx\csname date#1\endcsname\relax
  \else\selectlanguage{#1}\fi
}
\providecommand*{\BibUrl}[1]{\url{#1}}
\providecommand{\BibAnnote}[1]{}
\providecommand*{\BibEmph}[1]{#1}
\ProvideTextCommandDefault{\cyrdash}{\iflanguage{russian}{\hbox
  to.8em{--\hss--}}{\textemdash}}
\providecommand*{\BibDash}{}
\providecommand{\newblock}{\ignorespaces}
\makeatletter
\newcommand{\fullciteentry}[1]{%
  \@ifundefined{fullcite@#1}{%
    \textbf{[Источник #1 не найден; обновите полные библиографические записи]}%
  }{%
    \csname fullcite@#1\endcsname
  }%
}
\newcommand{\printfootcitelist}[1]{%
  \begingroup
  \def\href##1##2{##2}%
  \def\url##1{\nolinkurl{##1}}%
  \def\BibUrl##1{\nolinkurl{##1}}%
  \def\footcite@separator{}%
  \@for\footcite@key:=#1\do{%
    \footcite@separator\expandafter\fullciteentry\expandafter{\footcite@key}%
    \def\footcite@separator{\par\medskip}%
  }%
  \endgroup
}
\DeclareRobustCommand{\footcite}[1]{\footnote{\nocite{#1}\printfootcitelist{#1}}\selectlanguage{russian}}
\makeatother
\InputIfFileExists{biblio/footcite-entries.tex}{}{}

% Прочее
\usepackage{color} % работа с цветом
\usepackage{epstopdf} % конвертация eps в pdf
\usepackage{multirow} % объединение ячеек таблиц по вертикали
\usepackage{afterpage} % вставка материала после текущей страницы
\usepackage[font={normal}]{caption} % настройка подписей к рисункам и таблицам
\usepackage[onehalfspacing]{setspace} % полуторный интервал
\usepackage{fancyhdr} % установка колонтитулов
\usepackage{listings} % поддержка вставки исходного кода

% Установка шрифта Times New Roman
\renewcommand{\rmdefault}{ftm}

% Настройка текста
\usepackage{setspace}
\setstretch{1.5} % промежуточный интервал
\setlength{\parindent}{1.5cm}


% Создание нового типа столбца для выравнивания содержимого по центру
\newcommand{\PreserveBackslash}[1]{\let\temp=\\#1\let\\=\temp}
\newcolumntype{C}[1]{>{\PreserveBackslash\centering}p{#1}}
\newcolumntype{L}[1]{>{\raggedright\arraybackslash}p{#1}}
\newcolumntype{Y}{>{\raggedright\arraybackslash}X}

% Настройка стиля страницы
\pagestyle{fancy}      % Использование стиля "fancy" для оформления страниц
\fancyhf{}              % Очистка текущих значений колонтитулов
\fancyfoot[C]{\thepage} % Установка номера страницы в нижнем колонтитуле по центру
\renewcommand{\headrulewidth}{0pt} % Удаление разделительной линии в верхнем колонтитуле

% Настройка подписей к изображениям и таблицам
\captionsetup{format=hang,labelsep=period}
\DeclareCaptionLabelSeparator{tablelinebreak}{\par}
\captionsetup[table]{
  position=top,
  format=plain,
  labelsep=tablelinebreak,
  justification=centering,
  singlelinecheck=false
}

% Использование полужирного начертания для векторов
\let\vec=\mathbf

% Установка глубины оглавления
\setcounter{tocdepth}{2}

% Указание папки для поиска изображений
\graphicspath{{images/}}

% Установка стилей страницы и главы
\pagestyle{footcenter}
\chapterpagestyle{footcenter}


\begin{document}

\includepdf[pages={1-1}]{extra/Annotation\space Title.pdf}

\newpage
\renewcommand{\contentsname}{\centerline{\large СОДЕРЖАНИЕ}}
\selectlanguage{russian}
\tableofcontents

\section*{Введение}
\addcontentsline{toc}{section}{Введение}

Выпускная квалификационная работа посвящена реконструкции бинарных визуальных
стимулов $6 \times 6$ пикселей по многоканальным сигналам EEG,
зарегистрированным в фазе воспоминания ранее предъявленного изображения.
Постановка продолжает проект по реконструкции визуальных стимулов по
EEG~\footcite{dementyev2025visualstimuli}: готовый корпус используется для
строгой воспроизводимой оценки реконструкции, а не для повторного сбора данных.

Актуальность темы связана с развитием декодирования мозговой активности и
необходимостью проверять, содержит ли EEG-сигнал информацию о структуре
визуального образа. Из-за низкого отношения сигнал/шум и межсубъектной
вариативности такая реконструкция требует контроля утечек данных, сравнения с
простыми базовыми моделями и консервативной статистической
интерпретации~\footcite{10.1088/1741-2552/aab2f2}. Цель работы - разработать
воспроизводимый пайплайн реконструкции бинарных стимулов по EEG и проверить
его при контроле субъектной структуры данных.

\section*{Краткое описание основной части работы}
\addcontentsline{toc}{section}{Краткое описание основной части работы}

Основная часть работы построена от обзора предметной области к практической
реализации и экспериментальной проверке. В первой главе рассматриваются основы
EEG, различие между visual perception и visual imagery, а также современные
EEG-to-image подходы, включая EEG2IMAGE, DreamDiffusion, BrainVis и
GWIT~\footcite{10.48550/arxiv.2412.19999}. Сделан вывод, что для малых
imagery-датасетов особенно важны проверяемые постановки и протоколы с
контролем утечек.

Во второй главе задача сведена к восстановлению 36 бинарных пикселей. Описан
ранее собранный корпус: 16 испытуемых, 31 экспериментальная проба, 651
синхронизированное наблюдение и 217 случайных паттернов в исходном описании.
Основные эксперименты используют 180 наблюдений случайных стимулов фазы
воспоминания. Полное 15-секундное окно рассматривается как консервативная
единица анализа, согласованная с разметкой.

Для моделирования используются временные, спектральные, пространственные и
локальные признаки, спектральные представления FFT/Welch, Morlet, Superlet и
STFT, классические модели машинного обучения, логистическая регрессия как
референсная модель и компактные сверточные архитектуры.

Практическая часть оценивает качество в межсубъектном и двунаправленном
межпробном протоколах. Контролируются пересечения по испытуемым или пробам,
ключам наблюдений, случайным зёрнам и полным целевым изображениям; все
обучаемые преобразования настраиваются только по обучающей части. Результаты
показывают, что пайплайн воспроизводимо проверяет задачу, но статистически
надежного превосходства над логистической регрессией не выявлено. Лучшие
описательные значения сбалансированной точности составили 0.518382 и 0.512011
для двух протоколов соответственно.

\section*{Заключение}
\addcontentsline{toc}{section}{Заключение}

В работе реализован исследовательский цикл: постановка задачи, подготовка
входов, обучение моделей, контроль утечек данных и статистическое сравнение.
Главный вклад заключается в строгой проверке реконструкционного эффекта при
контроле субъектной структуры данных. Вывод консервативен: текущие данные и
модели не дают статистически надежного улучшения относительно логистической
регрессии. Дальнейшая работа требует расширения данных, внешней валидации и
адаптации к испытуемым.

\renewcommand{\bibsection}{%
  \section*{Список использованной литературы}
  \addcontentsline{toc}{section}{Список использованной литературы}
}
\bibliographystyle{biblio/gost2008n}
\bibliography{biblio/bibliography}

\end{document}
